\section*{摘要}
多溶剂条件下的药物溶解度受分子结构和实验条件共同影响。同一药物的重复测量及母体间采样不均，使未见药物母体上的误差需要按母体评价。本文从基础描述符 ExtraTrees 出发，比较丰富描述符、显式药物--溶剂交互和母体加权训练，并以 TabPFN 提供补充预测。两条分支的配置与融合权重在内层母体分组验证中选择。在 1507 条开发记录、128 个母体的嵌套评价中，母体等权 MAE 从基础树的 0.7469 降至丰富特征树的 0.6831，融合流程取得 0.6668。融合较丰富特征树的误差差值为 $-0.0163$，95\% 母体配对区间为 $[-0.0370,0.0001]$。低溶解度端的观测值容易被高估。本文结合树分支归因、温度响应和实测组合内关联解释模型行为，并将固定权重模型部署到交互应用。

\noindent\textbf{关键词}\quad 条件溶解度；ExtraTrees；TabPFN；母体分组验证；模型解释
